% Original: PDF strana 5, donji slajd.
\slajdid{02-s010}
\begin{frame}[label=02-s010]{Koraci sinteze kombinacione mreže}
% element: 02-s010:e01
Na osnovu zahteva za digitalni sistem:
\par\vspace{2mm}
\begin{enumerate}
\item Eventualno formirati funkcionalnu tabelu.
\item Eventualno napisati logičke funkcije.
\item Prikazati sistem električnom šemom, upotrebom simbola logičkih funkcija.
\end{enumerate}
\par\vspace{4mm}
\begin{columns}[T,onlytextwidth]
\begin{column}{.48\textwidth}
% element: 02-s010:e02
\Oznaka{Zašto „eventualno“?}\vspace{2mm}
Neki koraci se mogu preskočiti: iz tabele se može neposredno crtati šema, a ponekad i iz zahteva.
\end{column}
\begin{column}{.48\textwidth}
% element: 02-s010:e03
\Oznaka{Šema i stvarno kolo}\vspace{2mm}
Svakom simbolu logičke funkcije odgovara realno kolo sa električnim karakteristikama.
\end{column}
\end{columns}
\par\vspace{4mm}
% element: 02-s010:e04
\Rezultat{Međukorak: minimizacija ili prilagođavanje funkcije — broj tranzistora, kašnjenje, broj kola iste vrste.}
\end{frame}
